% ============================================================
% 答题卡模板
% 用于单独的答题卡，方便批改
% 编译命令: xelatex answer-sheet.tex
% ============================================================

\documentclass[11pt,a4paper]{article}

% ============ 页面设置 ============
\usepackage[top=2cm,bottom=2cm,left=2.5cm,right=2.5cm]{geometry}

% ============ 中文支持 ============
\usepackage{ctex}

% ============ 其他 ============
\usepackage{tikz}
\usepackage{array}
\usepackage{tabularx}
\usepackage{fancyhdr}
\usepackage{lastpage}
\usepackage{amssymb}
\usepackage{pifont}

% ============ 字体设置 ============
% macOS
\setCJKmainfont{Songti SC}[BoldFont={Heiti SC}]
\setCJKsansfont{Heiti SC}
\setmainfont{Times New Roman}

% ============ 页眉页脚 ============
\pagestyle{fancy}
\fancyhf{}
\renewcommand{\headrulewidth}{0.4pt}
\fancyhead[L]{\small XXX 公司 YYY 岗位笔试 — 答题卡}
\fancyhead[R]{\small\color{red}内部资料}
\fancyfoot[C]{\small 第 \thepage\ 页 / 共 \pageref{LastPage} 页}

% ============ 选项圆圈 ============
\newcommand{\optcircle}{\tikz[baseline=-0.5ex]\draw (0,0) circle (0.35em);}

% ============================================================
\begin{document}

% ============ 标题 ============
\begin{center}
{\LARGE\bfseries\sffamily XXX 公司 YYY 岗位笔试}

\vspace{0.3em}
{\Large\bfseries 答题卡}

\vspace{0.5em}
\begin{tabularx}{0.8\textwidth}{Xr}
姓名：\underline{\hspace{8em}} & 日期：\underline{\hspace{8em}}
\end{tabularx}
\end{center}

\vspace{0.5em}
\noindent\rule{\linewidth}{0.4pt}

% ============ 填涂说明 ============
\vspace{0.5em}
\noindent\textbf{注意事项：}
\begin{itemize}
\item 单选题：在对应选项上打 $\checkmark$ 或涂黑
\item 多选题：在所有正确选项上打 $\checkmark$ 或涂黑
\item 判断题：正确打 $\checkmark$，错误打 $\times$
\item 请使用黑色签字笔或 2B 铅笔填涂
\end{itemize}

\vspace{1em}

% ============ 一、单选题 ============
\noindent\textbf{一、单选题}（每题 3 分，共 30 分）

\vspace{0.5em}
\begin{tabular}{|c|c|c|c|c||c|c|c|c|c|}
\hline
\textbf{题号} & \textbf{A} & \textbf{B} & \textbf{C} & \textbf{D} &
\textbf{题号} & \textbf{A} & \textbf{B} & \textbf{C} & \textbf{D} \\
\hline
1 & \optcircle & \optcircle & \optcircle & \optcircle &
6 & \optcircle & \optcircle & \optcircle & \optcircle \\
\hline
2 & \optcircle & \optcircle & \optcircle & \optcircle &
7 & \optcircle & \optcircle & \optcircle & \optcircle \\
\hline
3 & \optcircle & \optcircle & \optcircle & \optcircle &
8 & \optcircle & \optcircle & \optcircle & \optcircle \\
\hline
4 & \optcircle & \optcircle & \optcircle & \optcircle &
9 & \optcircle & \optcircle & \optcircle & \optcircle \\
\hline
5 & \optcircle & \optcircle & \optcircle & \optcircle &
10 & \optcircle & \optcircle & \optcircle & \optcircle \\
\hline
\end{tabular}

\vspace{1em}

% ============ 二、多选题 ============
\noindent\textbf{二、多选题}（每题 4 分，共 20 分，少选错选多选不得分）

\vspace{0.5em}
\begin{tabular}{|c|c|c|c|c|}
\hline
\textbf{题号} & \textbf{A} & \textbf{B} & \textbf{C} & \textbf{D} \\
\hline
11 & \optcircle & \optcircle & \optcircle & \optcircle \\
\hline
12 & \optcircle & \optcircle & \optcircle & \optcircle \\
\hline
13 & \optcircle & \optcircle & \optcircle & \optcircle \\
\hline
14 & \optcircle & \optcircle & \optcircle & \optcircle \\
\hline
15 & \optcircle & \optcircle & \optcircle & \optcircle \\
\hline
\end{tabular}

\vspace{1em}

% ============ 三、判断题 ============
\noindent\textbf{三、判断题}（每题 2 分，共 20 分）

\vspace{0.5em}
\begin{tabular}{|c|c|c||c|c|c|}
\hline
\textbf{题号} & $\checkmark$ & $\times$ &
\textbf{题号} & $\checkmark$ & $\times$ \\
\hline
16 & \optcircle & \optcircle & 21 & \optcircle & \optcircle \\
\hline
17 & \optcircle & \optcircle & 22 & \optcircle & \optcircle \\
\hline
18 & \optcircle & \optcircle & 23 & \optcircle & \optcircle \\
\hline
19 & \optcircle & \optcircle & 24 & \optcircle & \optcircle \\
\hline
20 & \optcircle & \optcircle & 25 & \optcircle & \optcircle \\
\hline
\end{tabular}

\vspace{1em}

% ============ 四、简答题 ============
\noindent\textbf{四、简答题}（共 30 分）

\vspace{0.5em}
\noindent 请在试卷上作答，此处仅供阅卷登分使用。

\vspace{0.5em}
\begin{tabular}{|c|c|c|c|}
\hline
\textbf{题号} & \textbf{满分} & \textbf{得分} & \textbf{阅卷人} \\
\hline
26 & 10 & & \\
\hline
27 & 10 & & \\
\hline
28 & 10 & & \\
\hline
\end{tabular}

\vspace{2em}

% ============ 总分 ============
\noindent\rule{\linewidth}{0.4pt}

\vspace{0.5em}
\begin{center}
\begin{tabular}{|c|c|c|c|c||c|}
\hline
\textbf{单选题} & \textbf{多选题} & \textbf{判断题} & \textbf{简答题} & \textbf{总分} & \textbf{核分人} \\
\hline
/30 & /20 & /20 & /30 & /100 & \\
\hline
\end{tabular}
\end{center}

\end{document}
